\subsection{Programme and Path Governance}
\label{sec:impl_path_governance}

Section~\ref{sec:impl_career_paths} covers the pathway itself --- its versions, stages and the evaluation that decides what a student has finished. This section covers the enrolment structure above it: which career paths a student holds, how they are acquired, and how one is given up. It is implemented in \path{abridgeai/features/learning_programs/}. Figure~\ref{fig:path_change_attempt_closure} in Appendix~\ref{appendix:algorithm} traces the giving-up case in full, from the learner's request through the dean's decision to the entitlements the approval revokes.

\paragraph{A held path is a row, not a column}
The enrolment carries no \texttt{career\_path\_id}. \texttt{ProgramEnrollment} owns many \texttt{ProgramPathAttempt} rows, and the paths a student currently holds are the active ones. This is what lets a programme permit several paths at once, and it is what keeps a path that has ended readable: closing an attempt writes an \texttt{exit\_snapshot} of the progress it had earned rather than clearing a field --- measured by the same stage evaluation the learner's own view and the manager's roster report (Section~\ref{sec:pathway_completion}), so the frozen figure means what the live one meant the moment before it froze. Each attempt also records a \texttt{selection\_source} --- \texttt{program\_default}, \texttt{student}, or \texttt{path\_change} --- so an audit can distinguish a path the system assigned from one the student chose.

\paragraph{Automatic placement at enrolment}
A published programme version must nominate exactly one of its paths as the default, enforced by a partial unique index on \texttt{learning\_program\_version\_paths} rather than in application code, so no concurrent edit can produce a second one. Enrolment resolves that default \emph{before} writing any row: if it has since been archived, the whole operation aborts. Resolving it afterwards would have admitted the student and then failed, leaving an enrolment with no path --- a state nothing else in the feature can produce or recover from. Both enrolment doors, the hand-picked list and the CSV import, go through the same helper for this reason.

The default is left unstarted, rather than the enrolment refused, in three cases: the student already holds that path in another programme, they are at the organisation's ceiling for concurrent paths, or they have already completed the version the programme pins. The enrolment stands and waits in \texttt{awaiting\_path} so the student chooses a path for themselves. The last case is what stops a programme completing itself the moment it is joined: an attempt whose pinned version is already finished completes at the next progress evaluation, which would leave no active attempt, close the enrolment against a programme in which the student had done nothing, and refuse every further path they tried to add. It is judged by the same \texttt{is\_version\_complete\_for\_user} that evaluation uses, against the same pinned version, so the guard and the sweep cannot reach opposite answers.

\paragraph{Adding is direct; leaving is reviewed}
Taking on a further path is a direct write, bounded by two limits. A programme may cap how many of its paths one enrolment holds; the cap is optional and, where set, may not exceed the paths the programme bundles. Where it is unset the programme imposes no limit of its own, and what binds is the organisation-configurable ceiling on concurrent paths a student may hold across all their programmes --- the only limit that sees the total, since per-programme caps would otherwise sum, and a student in two programmes capped at one and two could hold three. Leaving one --- by switching to another path, or by dropping it outright --- creates a \texttt{PathChangeRequest} that a Faculty Dean must decide.

Both kinds share one table, discriminated by a \texttt{kind} column, because they share everything that matters: the review queue, the partial unique index enforcing one open request per enrolment, the finite switch allowance an approval spends, and the rejection vocabulary. They differ only in whether a destination exists, so the target columns are nullable under a \texttt{CHECK} constraint that requires one for a change and forbids one for a drop. A separate table would have duplicated the queue and the allowance, and made the single-open-request rule unenforceable by any one constraint.

\paragraph{One path, one programme at a time}
A career path may be bundled by several programmes, but a student may hold it in only one of them at a time. The rule is a partial unique index on \texttt{(student\_id, career\_path\_id)} across active attempts, which requires \texttt{student\_id} on the attempt itself: it otherwise lives on the enrolment, one table away, where no constraint can reach it. A denormalised key that drifts would be worse than none, so a composite foreign key ties \texttt{(program\_enrollment\_id, student\_id)} back to the enrolment it was copied from, which makes an attempt whose student disagrees with its enrolment's unrepresentable rather than merely unlikely. The service refuses a duplicate first and names the programme already holding it; the index is what holds when two requests in different enrolments race, which no per-enrolment constraint can see.

\paragraph{Invariants are re-checked at the decision, not trusted from the request}
Every precondition is verified twice: once when the student files, and again when the dean decides. The gap between the two is measured in days, and the enrolment can move within it. Approving a change re-confirms that the source attempt is still active and that the target has not since been archived, in which case the request is marked \texttt{invalidated} rather than failed. Approving a drop re-confirms that another path would remain in progress, because the student's other path may have been completed, switched out, or dropped by an earlier approval since filing.

\paragraph{Entitlements follow the paths, not the request}
Course access granted by a programme is released only when no remaining attempt still justifies it, resolved by counting the student's other active attempts on that path rather than assumed from the path being left. With several paths held at once, a course shared between two of them must survive one ending. Completions are held separately again, as \texttt{CourseCompletionAward} rows attached to the course, so a result already earned is never recomputed when the path beneath it changes.

\paragraph{Programme completion}
An enrolment completes only when no active attempt remains, so a student holding two paths completes the programme on finishing the second, not the first. Each attempt is judged against the path version it was pinned to, never the path's current head, so a revision published mid-enrolment cannot retroactively change what a student must finish.

Whether an attempt's pinned version has been finished is not decided here. It is asked of the career-path feature through \texttt{is\_version\_complete\_for\_user}, which runs the stage evaluator specified in Section~\ref{sec:pathway_completion}: every required course satisfied, and each stage's elective quota met. Re-deriving it here would invite a stricter reading than the one the student sees --- testing that every course attached to the version was satisfied, say --- and the two would part company exactly where the quota bites --- a pathway of two required courses and five electives asking for one would show the student $100\%$, complete their career enrolment, and leave the programme attempt open against four electives nobody owed. Neither side would surface the divergence, each being individually self-consistent. Completion is one rule with one implementation for the same reason the request table is one table: the moment a second copy exists, the two drift.

One distinction is preserved rather than flattened. Evaluated for the student's own pathway, a stage counts as complete if it is \emph{latched} or currently satisfied, so completion never regresses once earned. Evaluated for a programme attempt there is no latch to read: the attempt pins its own version, which need not be the one the student's career enrolment --- and its latches --- point at. The attempt is therefore judged live, against the version it actually promised.

A version presenting no stages at all returns not-complete in both paths. An unmeasurable pathway has not been shown to be finished, and what these writes produce is a durable academic record rather than a progress bar; publication guarantees at least one stage per path and one course per stage, so no pathway a student can hold reaches that branch.

\paragraph{The request lifecycle}
One lifecycle serves both kinds. A request is filed \texttt{pending} and may be acknowledged into \texttt{in\_progress}, which signals to the student that it is being considered but decides nothing --- both states are open, both occupy the one-request-per-enrolment slot, and both are decidable, so a dean may rule straight from the queue without acknowledging first. From either, the dean moves it to \texttt{approved} or \texttt{rejected}; the student may withdraw it to \texttt{cancelled} while it remains open; and the system itself may end it as \texttt{invalidated} where the re-checks above fail at the decision. That last transition is the reason the lifecycle has six states rather than five: a request overtaken by events is neither the student's fault nor the dean's refusal, and recording it as a rejection would misattribute both.

The two kinds diverge only at approval. Approving a change closes the outgoing attempt with a progress snapshot and opens a replacement pinned to the target path version; approving a drop closes the attempt and opens nothing. Everything before that point --- the queue, the allowance, the open-request rule, the rejection vocabulary --- is identical, which is what makes the shared table worth its nullable target columns.
